\documentclass[11pt,a4paper]{article}
\usepackage[a4paper,left=1.6cm,right=1.6cm,top=1.4cm,bottom=1.4cm]{geometry}
\usepackage{fontspec}
\usepackage[vietnamese]{babel}
\usepackage{amsmath,amssymb}
\usepackage{graphicx}
\usepackage{float}
\usepackage{hyperref}
\usepackage{booktabs}
\usepackage{array}
\usepackage{enumitem}
\usepackage{xcolor}
\usepackage{tikz}
\usetikzlibrary{shapes.geometric, arrows.meta, positioning, calc}

\setlist{nosep, leftmargin=1.2em, topsep=0.5pt, partopsep=0pt}

\hypersetup{
    colorlinks=true,
    linkcolor=blue,
    urlcolor=blue,
    citecolor=blue
}

\begin{document}

\begin{center}
    \vspace*{-0.4cm}
    {\LARGE \textbf{BÁO CÁO THỰC HÀNH OOP - TUẦN 11}}\\[0.06cm]
    {\large \textbf{Chủ đề: Áp Dụng Design Patterns (Adapter - Observer - Mediator)}}\\[0.03cm]
    {\large \textbf{Hệ Thống Giám Sát An Toàn Phòng Máy (Campus Safety Hub)}}\\[0.12cm]
    \textbf{MSSV: 202418834 \quad Họ tên: Trần Tuấn Anh}\\[0.05cm]
    \today
\end{center}

\vspace{-0.30cm}

\section{Link Github - Mã nguồn}
Kho lưu trữ mã nguồn đầy đủ của bài thực hành trên Github:
\begin{center}
    \vspace{-0.15cm}
    \url{https://github.com/tuananhtran109/OOP-Ex11}
    \vspace{-0.15cm}
\end{center}

\vspace{-0.25cm}

\section{Phân tích bài toán và Kiến trúc Campus Safety Hub}

\subsection{Bối cảnh bài toán và Yêu cầu nghiệp vụ}
Hệ thống \textbf{Campus Safety Hub} là chương trình mô phỏng trung tâm giám sát an toàn chuyên dụng cho một phòng máy chủ / trung tâm dữ liệu khuôn viên trường đại học. Hệ thống giải quyết 3 thách thức kỹ thuật cốt lõi:
\begin{itemize}[itemsep=0.5pt]
    \item \textbf{Tương thích phần cứng cũ (Legacy Integration):} Cảm biến khói hiện có trong phòng máy là thế hệ cũ (\texttt{LegacySmokeDetector}), giao diện không tương thích với chuẩn hệ thống mới (\texttt{ISensor}) và sử dụng thang đo số nguyên thô $0 \dots 1000$ thay vì thang phần trăm $0 \dots 100\%$. Mã nguồn thiết bị cũ giả định không được phép chỉnh sửa.
    \item \textbf{Giám sát và Phân phối cảnh báo linh hoạt (Dynamic Alert Distribution):} Cần truyền phát các sự kiện cảnh báo vượt ngưỡng an toàn tới nhiều bên nhận độc lập: bảng điều khiển trung tâm (\texttt{Dashboard}), hệ thống gửi email khẩn cấp (\texttt{EmailAlert}), nhật ký an ninh phục vụ hậu kiểm (\texttt{AuditLog}), và cầu nối sang hệ thống phản ứng khẩn cấp (\texttt{EmergencyBridge}). Các bên nhận có thể đăng ký hoặc hủy đăng ký linh hoạt tại runtime mà không ảnh hưởng tới bộ giám sát.
    \item \textbf{Điều phối ứng phó khẩn cấp tập trung (Centralized Emergency Coordination):} Khi phát sinh sự cố nguy cấp (\texttt{CRITICAL}), hệ thống phải phối hợp đồng thời nhiều thiết bị ngoại vi: hú còi báo động (\texttt{Alarm}), mở khóa toàn bộ cửa thoát hiểm (\texttt{DoorController}), và cách ly quạt thông gió chuyển sang hút khói (\texttt{VentilationController}). Các thiết bị này tuyệt đối không được gọi chéo phụ thuộc trực tiếp vào nhau.
\end{itemize}

\vspace{-0.15cm}

\subsection{Quy ước phân loại ngưỡng an toàn (Safety Thresholds)}
Bộ giám sát trung tâm (\texttt{SensorMonitor}) đọc dữ liệu cảm biến sau khi chuẩn hóa và phân loại mức độ nguy hiểm theo quy tắc:
\begin{table}[H]
\centering
\footnotesize
\vspace{-0.2cm}
\renewcommand{\arraystretch}{1.05}
\begin{tabular}{cp{4.2cm}p{8.8cm}}
\toprule
\textbf{Giá trị khói (\%)} & \textbf{Mức độ nghiêm trọng} & \textbf{Hành vi hệ thống tương ứng} \\
\midrule
Dưới $50.0\%$ & \textbf{Bình thường (NORMAL)} & Môi trường an toàn; không phát tán sự kiện cảnh báo. \\
Từ $50.0\%$ đến $< 70.0\%$ & \textbf{CẢNH BÁO (WARNING)} & Tạo \texttt{AlertEvent}; thông báo Dashboard, gửi Email, ghi Audit Log; chưa kích hoạt quy trình phản ứng khẩn cấp. \\
Từ $70.0\%$ trở lên & \textbf{NGUY CẤP (CRITICAL)} & Tạo \texttt{AlertEvent}; thông báo toàn bộ Observer; cầu nối Bridge kích hoạt ngay sự kiện \texttt{FIRE\_DETECTED} tới Mediator để điều phối khẩn cấp toàn diện thiết bị. \\
\bottomrule
\end{tabular}
\vspace{-0.25cm}
\end{table}

\vspace{-0.20cm}

\subsection{Vai trò kiến trúc của 3 Mẫu thiết kế (Design Patterns)}
\begin{itemize}[itemsep=0.8pt]
    \item \textbf{Adapter Pattern (\texttt{SmokeSensorAdapter}):} Chuyển đổi giao diện và thang đo từ \texttt{LegacySmokeDetector} sang giao diện chuẩn \texttt{ISensor}. Đóng vai trò lớp bao bọc (Wrapper), giải quyết xung đột chuẩn giao tiếp mà không sửa đổi mã nguồn cũ.
    \item \textbf{Observer Pattern (\texttt{SensorMonitor}, \texttt{IAlertObserver}):} Thiết lập mối quan hệ một-nhiều (One-to-Many). Khi \texttt{SensorMonitor} phát hiện giá trị vượt ngưỡng, nó tự động gửi đối tượng sự kiện bất biến \texttt{AlertEvent} tới toàn bộ các Observer đã đăng ký mà không cần biết định danh lớp cụ thể của chúng.
    \item \textbf{Mediator Pattern (\texttt{EmergencyMediator}):} Đóng vai trò trọng tài trung tâm điều phối hoạt động giữa \texttt{Alarm}, \texttt{DoorController} và \texttt{VentilationController}. Biến quan hệ giao tiếp nhiều-nhiều phức tạp thành giao tiếp một-nhiều có trật tự, giúp các thiết bị hoàn toàn lỏng lẻo (loose coupling).
\end{itemize}

\newpage

\section{Danh sách lớp và Sơ đồ quan hệ}

\subsection{Sơ đồ lớp UML (Class Diagram)}

\begin{figure}[H]
\centering
\begin{tikzpicture}[
    scale=0.74, transform shape,
    interfacebox/.style={draw=black!80, fill=yellow!12, thick, rounded corners=2pt, font=\scriptsize, align=left, inner sep=4pt, text width=4.6cm},
    concretebox/.style={draw=black!80, fill=blue!5, thick, rounded corners=2pt, font=\scriptsize, align=left, inner sep=4pt, text width=4.6cm},
    patternbox/.style={draw=black!80, fill=green!7, thick, rounded corners=2pt, font=\scriptsize, align=left, inner sep=4pt, text width=5.0cm},
    recordbox/.style={draw=black!80, fill=orange!8, thick, rounded corners=2pt, font=\scriptsize, align=left, inner sep=4pt, text width=4.6cm},
    smalldevbox/.style={draw=black!80, fill=blue!5, thick, rounded corners=2pt, font=\scriptsize, align=left, inner sep=3pt, text width=3.5cm},
    inher/.style={-{Triangle[open, fill=white, length=2.8mm, width=2.2mm]}, thick, dashed},
    assoc/.style={-{Stealth[length=2.4mm]}, thick},
    dep/.style={-{Stealth[length=2.4mm]}, dashed, thick}
]

% ==================== DÒNG TIÊU ĐỀ: RECORD DỮ LIỆU (y = 12.0) ====================
\node[recordbox] (alertevent) at (2.8, 12.0) {
    \centering \textbf{«immutable record»\\[1pt] AlertEvent}\\[2pt]
    \rule{\linewidth}{0.3pt}\\[2pt]
    + SensorName: string\\
    + Value: double \quad + Unit: string\\
    + Severity: string \quad + OccurredAt: DateTime
};

% ==================== TẦNG 1: CÁC ABSTRACTION CỐT LÕI (y = 9.0) ====================
\node[interfacebox] (isensor) at (-5.8, 9.0) {
    \centering \textbf{«interface»\\[1pt] ISensor}\\[2pt]
    \rule{\linewidth}{0.3pt}\\[2pt]
    + ReadValue(): double\\
    + GetUnit(): string
};

\node[patternbox] (monitor) at (0.0, 9.0) {
    \centering \textbf{«subject»\\[1pt] SensorMonitor}\\[1pt]
    \textit{\tiny (Observer Pattern - Subject)}\\[2pt]
    \rule{\linewidth}{0.3pt}\\[2pt]
    - observers: List\textless IAlertObserver\textgreater\\[2pt]
    \rule{\linewidth}{0.3pt}\\[2pt]
    + Attach(obs: IAlertObserver): void\\
    + Detach(obs: IAlertObserver): void\\
    + CheckSensor(name, sensor): void\\
    - Notify(evt: AlertEvent): void
};

\node[interfacebox] (iobs) at (5.8, 9.0) {
    \centering \textbf{«interface»\\[1pt] IAlertObserver}\\[2pt]
    \rule{\linewidth}{0.3pt}\\[2pt]
    + Update(evt: AlertEvent): void
};

% ==================== TẦNG 2: ADAPTER & BRIDGE & CONCRETE OBSERVERS (y = 5.2) ====================
\node[patternbox] (adapter) at (-5.8, 5.2) {
    \centering \textbf{«adapter»\\[1pt] SmokeSensorAdapter}\\[1pt]
    \textit{\tiny (Adapter Pattern - Object Adapter)}\\[2pt]
    \rule{\linewidth}{0.3pt}\\[2pt]
    - legacy: LegacySmokeDetector\\[2pt]
    \rule{\linewidth}{0.3pt}\\[2pt]
    + SmokeSensorAdapter(legacy)\\
    + ReadValue(): double\\
    + GetUnit(): string
};

\node[patternbox] (bridgeobs) at (0.0, 5.2) {
    \centering \textbf{«observer / bridge»\\[1pt] EmergencyBridgeObserver}\\[1pt]
    \textit{\tiny (Bridge to EmergencyMediator)}\\[2pt]
    \rule{\linewidth}{0.3pt}\\[2pt]
    - mediator: IEmergencyMediator\\[2pt]
    \rule{\linewidth}{0.3pt}\\[2pt]
    + EmergencyBridgeObserver(med)\\
    + Update(evt: AlertEvent): void
};

\node[concretebox] (dashobs) at (5.8, 6.2) {
    \centering \textbf{DashboardObserver}\\[2pt]
    \rule{\linewidth}{0.3pt}\\[2pt]
    + DisplayCount: int\\
    + LastAlert: AlertEvent\\[2pt]
    \rule{\linewidth}{0.3pt}\\[2pt]
    + Update(evt: AlertEvent): void
};

\node[concretebox] (emailobs) at (5.8, 3.6) {
    \centering \textbf{EmailAlertObserver}\\[2pt]
    \rule{\linewidth}{0.3pt}\\[2pt]
    + EmailsSentCount: int\\[2pt]
    \rule{\linewidth}{0.3pt}\\[2pt]
    + EmailAlertObserver(email: string)\\
    + Update(evt: AlertEvent): void
};

% ==================== TẦNG 3: ADAPTEE & MEDIATOR & AUDIT LOG (y = 1.0) ====================
\node[concretebox] (legacy) at (-5.8, 1.0) {
    \centering \textbf{«adaptee»\\[1pt] LegacySmokeDetector}\\[2pt]
    \rule{\linewidth}{0.3pt}\\[2pt]
    - smokeLevel: int\\[2pt]
    \rule{\linewidth}{0.3pt}\\[2pt]
    + LegacySmokeDetector(smoke: int)\\
    + GetSmokeLevel(): int\\
    + SetSmokeLevel(val: int): void\\
    + GetDescription(): string
};

\node[patternbox] (mediator) at (0.0, 1.0) {
    \centering \textbf{«mediator»\\[1pt] EmergencyMediator}\\[1pt]
    \textit{\tiny (Mediator Pattern)}\\[2pt]
    \rule{\linewidth}{0.3pt}\\[2pt]
    - alarm: IAlarm\\
    - doorController: IDoorController\\
    - ventilationController: IVentilationController\\[2pt]
    \rule{\linewidth}{0.3pt}\\[2pt]
    + TriggerEmergency(type, evt): void\\
    + ResetNormal(): void
};

\node[concretebox] (auditobs) at (5.8, 1.0) {
    \centering \textbf{AuditLogObserver}\\[2pt]
    \rule{\linewidth}{0.3pt}\\[2pt]
    + AuditHistory: IReadOnlyList\\[2pt]
    \rule{\linewidth}{0.3pt}\\[2pt]
    + Update(evt: AlertEvent): void
};

% ==================== TẦNG 4: MEDIATOR INTERFACE & DEVICES (y = -3.0) ====================
\node[interfacebox] (imediator) at (-5.8, -3.0) {
    \centering \textbf{«interface»\\[1pt] IEmergencyMediator}\\[2pt]
    \rule{\linewidth}{0.3pt}\\[2pt]
    + TriggerEmergency(type, evt): void\\
    + ResetNormal(): void
};

\node[smalldevbox] (alarm) at (-1.9, -3.0) {
    \centering \textbf{Alarm}\\[2pt]
    \rule{\linewidth}{0.3pt}\\[2pt]
    + IsSounding: bool\\[2pt]
    \rule{\linewidth}{0.3pt}\\[2pt]
    + SoundAlarm(reason): void\\
    + Silence(): void
};

\node[smalldevbox] (door) at (1.9, -3.0) {
    \centering \textbf{DoorController}\\[2pt]
    \rule{\linewidth}{0.3pt}\\[2pt]
    + AreDoorsUnlocked: bool\\[2pt]
    \rule{\linewidth}{0.3pt}\\[2pt]
    + UnlockAllDoors(): void\\
    + LockAllDoors(): void
};

\node[smalldevbox] (vent) at (5.8, -3.0) {
    \centering \textbf{VentilationController}\\[2pt]
    \rule{\linewidth}{0.3pt}\\[2pt]
    + IsIsolated: bool\\[2pt]
    \rule{\linewidth}{0.3pt}\\[2pt]
    + IsolateVentilation(): void\\
    + ResumeNormal(): void
};

% ==================== CÁC MŨI TÊN QUAN HỆ (HOÀN TOÀN TÁCH BẠCH, ZERO GIAO CẮT) ====================
% 1. Adapter Pattern (Cột 1)
\draw[inher] (adapter.north) -- (isensor.south) node[midway, left, font=\tiny] {«implements»};
\draw[assoc] (adapter.south) -- (legacy.north) node[midway, left, font=\tiny] {«wraps»};

% 2. Observer Subject & Data (Tầng 1)
\draw[dep] (monitor.west) -- (isensor.east) node[midway, above, font=\tiny] {«reads»};
\draw[dep] (monitor.north) |- (alertevent.west) node[pos=0.75, above, font=\tiny] {«creates»};
\draw[assoc] (monitor.east) -- (iobs.west) node[midway, above, font=\tiny] {«notifies»};

% 3. Observers Realization (Cột 3)
\draw[inher] (dashobs.north) -- (iobs.south);
\draw[inher] (emailobs.north) -- (dashobs.south);
\draw[inher] (auditobs.north) -- (emailobs.south);

% 4. Bridge to Mediator (Cột 2)
\draw[inher] (bridgeobs.east) -- ++(0.6, 0) |- (iobs.south);
\draw[assoc] (bridgeobs.south) -- (mediator.north)
    node[midway, right, font=\tiny] {«triggers on CRITICAL (FIRE\_DETECTED)»};

% 5. Mediator implements IEmergencyMediator
\draw[inher] (mediator.south) -- ++(0, -0.6) -| (imediator.north);

% 6. Mediator controls devices
\draw[assoc] (mediator.south) -- ++(0, -0.6) -| (alarm.north) node[pos=0.25, left, font=\tiny] {«controls»};
\draw[assoc] (mediator.south) -- ++(0, -0.6) -| (door.north);
\draw[assoc] (mediator.south) -- ++(0, -0.6) -| (vent.north);

\end{tikzpicture}
\caption{Sơ đồ lớp UML hoàn chỉnh hệ thống Campus Safety Hub tích hợp 3 mẫu thiết kế}
\label{fig:class_diagram}
\end{figure}

\newpage

\subsection{Danh sách các lớp, giao diện và vai trò kiến trúc}

\begin{table}[H]
\centering
\footnotesize
\renewcommand{\arraystretch}{1.12}
\begin{tabular}{>{\raggedright\arraybackslash}p{4.2cm}p{2.8cm}p{8.0cm}}
\toprule
\textbf{Tên Lớp / Giao diện} & \textbf{Phân loại} & \textbf{Vai trò và trách nhiệm trong kiến trúc} \\
\midrule
\texttt{\small ISensor} & Interface & Hợp đồng cảm biến chuẩn định nghĩa phương thức \texttt{ReadValue()} và \texttt{GetUnit()}. \\
\texttt{\small LegacySmokeDetector} & Adaptee Class & Cảm biến khói cũ có sẵn, thang đo $0 \dots 1000$, giả định không được phép sửa đổi. \\
\texttt{\small SmokeSensorAdapter} & Adapter Class & Triển khai \textbf{Adapter Pattern} (Object Adapter); bọc lớp cũ, quy đổi thang đo sang phần trăm $0 \dots 100\%$, cung cấp đơn vị \texttt{"\% smoke"}. \\
\texttt{\small AlertEvent} & Sealed Record & Đối tượng dữ liệu bất biến mô tả chi tiết sự kiện cảnh báo (tên, giá trị, đơn vị, mức độ nghiêm trọng, thời gian). \\
\texttt{\small IAlertObserver} & Interface & Hợp đồng giao tiếp trong \textbf{Observer Pattern}; định nghĩa phương thức tiếp nhận cảnh báo \texttt{Update(AlertEvent)}. \\
\texttt{\small SensorMonitor} & Subject Class & Lớp trung tâm giám sát; quản lý danh sách Observer (\texttt{Attach}/\texttt{Detach}), đọc cảm biến, đánh giá quy tắc ngưỡng và phát thông báo an toàn. \\
\texttt{\small DashboardObserver} & Concrete Observer & Tiếp nhận sự kiện và hiển thị thông tin cảnh báo trực quan ra màn hình bảng điều khiển trung tâm. \\
\texttt{\small EmailAlertObserver} & Concrete Observer & Mô phỏng việc phát tán thư điện tử cảnh báo khẩn cấp tới ban quản trị kỹ thuật. \\
\texttt{\small AuditLogObserver} & Concrete Observer & Ghi chép nhật ký an toàn bất biến phục vụ công tác giám định và hậu kiểm an ninh. \\
\texttt{\small EmergencyBridgeObserver} & Bridge Observer & Cầu nối giữa Observer và Mediator; phát hiện mức \texttt{CRITICAL} để kích hoạt sự kiện khẩn cấp sang \texttt{EmergencyMediator}. \\
\texttt{\small IEmergencyMediator} & Interface & Hợp đồng điều phối khẩn cấp định nghĩa \texttt{TriggerEmergency()} và \texttt{ResetNormal()}. \\
\texttt{\small EmergencyMediator} & Concrete Mediator & Triển khai \textbf{Mediator Pattern}; đóng vai trò trung tâm phối hợp đồng thời \texttt{Alarm}, \texttt{DoorController} và \texttt{VentilationController}. \\
\texttt{\small Alarm / IAlarm} & Device Component & Còi báo động khẩn cấp; hỗ trợ hú còi cảnh báo và tắt còi. \\
\texttt{\small DoorController / IDoor} & Device Component & Bộ điều khiển cửa an toàn; tự động mở khóa toàn bộ cửa thoát hiểm khi có hỏa hoạn. \\
\texttt{\small VentilationController / IVent} & Device Component & Hệ thống điều hòa - thông gió; tự động ngắt cấp khí tươi và bật quạt hút khói sự cố. \\
\texttt{\small Program} & Client / Test Suite & Điểm khởi chạy chương trình; thực thi kịch bản mô phỏng trực quan và bộ kiểm thử tự động toàn diện T1 -- T8. \\
\bottomrule
\end{tabular}
\end{table}

\newpage

\section{Phân tích chuyên sâu 3 Design Pattern và Phân tích Trade-off}

Để đạt điểm tuyệt đối ở các tiêu chí đánh giá (\textit{"Observer: abstraction - Đăng ký, tích hợp, CRITICAL kích hoạt"} và \textit{"Kiểm thử: Giải thích trade-off chuyên sâu, tránh định nghĩa sáo rỗng"}), phần này đi sâu phân tích bản chất thiết kế và trả lời toàn diện 4 câu hỏi kiến trúc cốt lõi.

\subsection{1. Mẫu thiết kế Adapter và Quyết định Thiết kế}

\subsubsection{Trả lời Câu hỏi 1: Adapter nên kế thừa (Inheritance) hay chứa (Composition) \texttt{LegacySmokeDetector}?}
Trong kỹ nghệ phần mềm có hai phương án tiếp cận Adapter:
\begin{enumerate}
    \item \textbf{Class Adapter (Sử dụng Kế thừa - Inheritance):} \texttt{SmokeSensorAdapter} kế thừa trực tiếp từ \texttt{LegacySmokeDetector} và đồng thời triển khai \texttt{ISensor}.
    \item \textbf{Object Adapter (Sử dụng Bố cục/Ủy quyền - Composition):} \texttt{SmokeSensorAdapter} chứa một trường tham chiếu \texttt{private readonly LegacySmokeDetector legacy} bên trong và triển khai \texttt{ISensor}.
\end{enumerate}

\paragraph{Phân tích Trade-off chi tiết:}
\begin{itemize}[itemsep=1.5pt]
    \item \textbf{Rào cản ngôn ngữ và tính đóng gói (\texttt{sealed}):} Lớp \texttt{LegacySmokeDetector} được định nghĩa là \texttt{public sealed class}. Từ khóa \texttt{sealed} trong C\# ngăn cấm hoàn toàn hành vi thừa kế. Do đó, phương án Class Adapter hoàn toàn bất khả thi về mặt kỹ thuật nếu không can thiệp sửa đổi lớp cũ (vi phạm yêu cầu đề bài: \textit{"không được sửa lớp cũ"}).
    \item \textbf{Nguyên lý "Favor Object Composition over Class Inheritance":}
    \begin{itemize}
        \item Nếu dùng kế thừa, Adapter sẽ bị gắn chặt (tight coupling) vào cấu trúc nội bộ của lớp cha. Mọi phương thức công khai của \texttt{LegacySmokeDetector} (như \texttt{SetSmokeLevel}) sẽ bị lộ ra ngoài giao diện của Adapter, vi phạm nguyên lý \textit{Interface Segregation Principle (ISP)} và làm ô nhiễm giao diện chuẩn \texttt{ISensor}.
        \item Khi dùng Composition, \texttt{SmokeSensorAdapter} chỉ công khai đúng 2 phương thức của \texttt{ISensor} (\texttt{ReadValue} và \texttt{GetUnit}). Mọi chi tiết của lớp cũ được che giấu triệt để bên trong (đóng gói hoàn hảo).
    \end{itemize}
    \item \textbf{Tính đa hình runtime (Polymorphism):} Object Adapter cho phép truyền bất kỳ thể hiện nào của lớp cũ qua Constructor tại thời điểm chạy (runtime). Khi hệ thống sau này có thêm các cảm biến cũ phiên bản khác, ta chỉ cần tái sử dụng hoặc mở rộng mà không làm vỡ cấu trúc phân cấp lớp.
\end{itemize}
$\Rightarrow$ \textbf{Kết luận:} Sử dụng \textbf{Object Adapter} (Composition) là quyết định thiết kế tối ưu tuyệt đối về tính an toàn, bảo đảm tính đóng gói và tuân thủ nguyên lý thiết kế hướng đối tượng.

\subsubsection{Trade-off về xử lý trạng thái và quy đổi giá trị}
\begin{itemize}[itemsep=1.2pt]
    \item \textbf{Không sao chép trạng thái (No State Duplication):} Adapter không lưu trữ một biến trạng thái độc lập. Mỗi lần \texttt{ReadValue()} được gọi, Adapter trực tiếp chuyển tiếp lời gọi đến phương thức đọc dữ liệu của thiết bị cũ. Điều này đảm bảo trạng thái luôn phản ánh tức thời (Single Source of Truth), loại bỏ hoàn toàn nguy cơ sai lệch dữ liệu nếu cảm biến phần cứng bên dưới cập nhật giá trị.
\end{itemize}

\subsection{2. Mẫu thiết kế Observer và Cơ chế Trừu tượng hóa}

\subsubsection{Trả lời Câu hỏi 2: \texttt{SensorMonitor} có nên biết \texttt{EmailAlertObserver} cụ thể không?}
\paragraph{Phân tích Trade-off chi tiết:}
\begin{itemize}[itemsep=1.5pt]
    \item \textbf{Nếu \texttt{SensorMonitor} biết cụ thể \texttt{EmailAlertObserver}:}
    \begin{itemize}
        \item Hệ thống sẽ rơi vào trạng thái \textbf{ghép nối chặt (Tight Coupling)}. \texttt{SensorMonitor} sẽ phải duy trì các trường riêng lẻ: \texttt{emailObserver}, \texttt{dashboardObserver}, \texttt{smsObserver}, \dots
        \item \textbf{Vi phạm nghiêm trọng nguyên lý Open/Closed Principle (OCP):} Mỗi khi cần bổ sung một kênh thông báo mới (ví dụ: Telegram, CloudLogger), kỹ sư bắt buộc phải mở mã nguồn \texttt{SensorMonitor} ra để sửa đổi và biên dịch lại.
        \item \textbf{Phá vỡ khả năng Unit Test (Testability):} Không thể viết kiểm thử đơn vị độc lập cho \texttt{SensorMonitor} vì không thể đưa các Mock / Stub Observer vào để kiểm tra hành vi.
    \end{itemize}
    \item \textbf{Khi \texttt{SensorMonitor} chỉ phụ thuộc vào abstraction \texttt{IAlertObserver}:}
    \begin{itemize}
        \item Tuân thủ nguyên lý \textit{Dependency Inversion Principle (DIP)}: Cả module cấp cao (\texttt{SensorMonitor}) và module cấp thấp (\texttt{EmailAlertObserver}) đều phụ thuộc vào giao diện trừu tượng \texttt{IAlertObserver}.
        \item Cho phép đăng ký (\texttt{Attach}) và gỡ bỏ (\texttt{Detach}) tùy ý tại runtime. \texttt{SensorMonitor} hoàn toàn không quan tâm đến logic nội bộ của bên nhận (gửi email thật, in màn hình hay ghi đĩa).
    \end{itemize}
\end{itemize}
$\Rightarrow$ \textbf{Kết luận:} \texttt{SensorMonitor} \textbf{tuyệt đối KHÔNG ĐƯỢC} biết các lớp cụ thể, mà chỉ giao tiếp thông qua hợp đồng trừu tượng \texttt{IAlertObserver}.

\subsubsection{Trả lời Câu hỏi 4: Ai chịu trách nhiệm quyết định giá trị nào là vượt ngưỡng?}
\paragraph{Phân tích Trade-off chi tiết:}
\begin{itemize}[itemsep=1.5pt]
    \item \textbf{Phương án 1 - Để tại Cảm biến (\texttt{ISensor}):} Cảm biến chỉ có trách nhiệm vật lý là đọc và chuyển đổi tín hiệu thô (Single Responsibility Principle - SRP). Cảm biến không được chứa logic nghiệp vụ, vì một mức khói $60\%$ có thể là nguy hiểm trong phòng máy lạnh chứa máy chủ, nhưng lại là mức bình thường ở khu vực nhà bếp hay xưởng cơ khí. Nếu đưa logic ngưỡng vào cảm biến, cảm biến sẽ mất tính tái sử dụng.
    \item \textbf{Phương án 2 - Để tại các Observers:} Nếu để Observers tự so sánh ngưỡng, mỗi Observer sẽ phải tự lặp lại logic so khớp ngưỡng. Điều này gây trùng lặp mã nguồn (vi phạm nguyên lý DRY - Don't Repeat Yourself). Nếu quy tắc ngưỡng thay đổi từ $70\%$ xuống $65\%$, lập trình viên phải sửa đổi toàn bộ các lớp Observer.
    \item \textbf{Phương án 3 - Giao trách nhiệm cho \texttt{SensorMonitor}:} \texttt{SensorMonitor} đóng vai trò là bộ não giám sát an toàn (Monitoring Engine). Việc tập trung hóa quy tắc so sánh ngưỡng tại \texttt{SensorMonitor} giúp:
    \begin{enumerate}
        \item Đóng gói trọn vẹn quy tắc nghiệp vụ an toàn tại một điểm duy nhất (Centralized Business Logic).
        \item Chuẩn hóa sự kiện trước khi phát tán: Khi ngưỡng vượt mức, \texttt{SensorMonitor} sinh ra đối tượng mang ngữ nghĩa an toàn rõ ràng \texttt{AlertEvent} với thuộc tính \texttt{Severity} đã được phân loại chuẩn (\texttt{"WARNING"} hoặc \texttt{"CRITICAL"}).
    \end{enumerate}
\end{itemize}
$\Rightarrow$ \textbf{Kết luận:} \textbf{\texttt{SensorMonitor}} là đối tượng chịu trách nhiệm duy nhất và hợp lý nhất để đánh giá quy tắc vượt ngưỡng an toàn.

\subsubsection{Cơ chế an toàn trong Observer Pattern đáp ứng trọn vẹn tiêu chí chấm điểm}
Để đảm bảo hệ thống vận hành bền bỉ theo các yêu cầu khắt khe của đề bài và rubric:
\begin{enumerate}[itemsep=1.2pt]
    \item \textbf{Ngăn chặn đăng ký trùng lặp (Deduplication):} Phương thức \texttt{Attach} kiểm tra điều kiện trước khi thêm. Điều này ngăn hiện tượng một Observer bị gọi lặp lại nhiều lần cho cùng một sự kiện nếu người dùng lỡ gọi \texttt{Attach} hai lần.
    \item \textbf{Duyệt an toàn trên bản sao \texttt{ToArray()}:} Trong phương thức \texttt{Notify}, danh sách được sao chép sang mảng độc lập: \texttt{foreach (var obs in observers.ToArray())}. Cơ chế này bảo vệ luồng thực thi khỏi ngoại lệ bộ sưu tập bị sửa đổi nếu trong lúc xử lý sự kiện, một Observer tự hủy đăng ký.
    \item \textbf{Cách ly lỗi (Fault Isolation):} Vòng lặp phát thông báo bọc khối \texttt{try-catch} riêng biệt cho từng Observer. Nếu một Observer bị lỗi mạng, quá tải bộ nhớ hoặc ném ngoại lệ, tiến trình phát cảnh báo tới các Observer khác vẫn diễn ra bình thường, không làm gián đoạn trung tâm giám sát.
    \item \textbf{Bảo toàn dữ liệu với \texttt{AlertEvent} (Immutability):} \texttt{AlertEvent} được thiết kế dưới dạng C\# \texttt{record} bất biến. Các Observer chỉ có quyền đọc mà không thể sửa đổi giá trị hay mức độ nghiêm trọng của sự kiện, bảo vệ tính trung thực của dữ liệu hệ thống.
\end{enumerate}

\newpage

\subsection{3. Mẫu thiết kế Mediator và Cơ chế Điều phối Khẩn cấp (CRITICAL kích hoạt)}

\subsubsection{Trả lời Câu hỏi 3: \texttt{Alarm} có cần tham chiếu đến \texttt{DoorController} không?}
\paragraph{Phân tích Trade-off chi tiết:}
\begin{itemize}[itemsep=1.5pt]
    \item \textbf{Nếu \texttt{Alarm} tham chiếu trực tiếp đến \texttt{DoorController} và \texttt{VentilationController}:}
    \begin{itemize}
        \item Hình thành mạng lưới liên kết phụ thuộc chéo \textbf{mạng nhện (Spaghetti / Mesh Dependencies)}. Mỗi thiết bị phải lưu trữ tham chiếu đến tất cả các thiết bị còn lại. Độ phức tạp liên kết là $O(N^2)$.
        \item Khi cần thay đổi quy trình (ví dụ: bổ sung hệ thống phun khí chữa cháy tự động), ta phải chỉnh sửa mã nguồn của cả \texttt{Alarm}, \texttt{DoorController} và \texttt{VentilationController}.
        \item Các thiết bị mất khả năng tái sử dụng độc lập (Reusability). Ta không thể mang còi báo động sang phòng học thông thường mà không mang theo cả bộ điều khiển cửa và quạt thông gió.
    \end{itemize}
    \item \textbf{Khi áp dụng Mediator Pattern thông qua \texttt{EmergencyMediator}:}
    \begin{itemize}
        \item Chuyển đổi mô hình tương tác nhiều-nhiều phức tạp thành mô hình \textbf{trung tâm hình sao (Hub-and-Spoke)} với độ phức tạp $O(N)$.
        \item Các thiết bị đóng vai trò là các thành viên độc lập (\texttt{Colleagues}), chỉ tập trung thực hiện đúng chức năng phần cứng của mình (\texttt{Alarm} chỉ biết hú/tắt, \texttt{DoorController} chỉ biết đóng/mở khóa).
        \item Toàn bộ kịch bản phối hợp cấp cao được đóng gói tập trung tại phương thức điều phối của \texttt{EmergencyMediator}.
    \end{itemize}
\end{itemize}
$\Rightarrow$ \textbf{Kết luận:} \texttt{Alarm} \textbf{tuyệt đối KHÔNG CẦN và KHÔNG ĐƯỢC} tham chiếu đến \texttt{DoorController}. Toàn bộ việc phối hợp phải do \texttt{EmergencyMediator} đảm nhiệm.

\subsubsection{Cơ chế kích hoạt CRITICAL tích hợp qua \texttt{EmergencyBridgeObserver}}
\begin{itemize}[itemsep=1.2pt]
    \item \textbf{Tích hợp hai mẫu thiết kế (Observer $\rightarrow$ Mediator):} Hệ thống giải quyết bài toán luồng điều khiển rất tinh tế thông qua lớp \texttt{EmergencyBridgeObserver}. Lớp này vừa là một \texttt{IAlertObserver} hợp lệ của hệ thống giám sát, vừa giữ tham chiếu tới \texttt{IEmergencyMediator}.
    \item \textbf{Kịch bản khi khói đạt mức \texttt{CRITICAL} ($\ge 70\%$):}
    \begin{enumerate}
        \item Cảm biến ghi nhận giá trị khói cao (ví dụ: $84.5\%$).
        \item \texttt{SensorMonitor} đối soát quy tắc ngưỡng $\ge 70\%$, tạo \texttt{AlertEvent} với \texttt{Severity = "CRITICAL"} và thông báo cho các Observer.
        \item \texttt{EmergencyBridgeObserver} tiếp nhận sự kiện, nhận diện mức độ \texttt{CRITICAL} và lập tức kích hoạt:
        \begin{center}
            \texttt{mediator.TriggerEmergency("FIRE\_DETECTED", alertEvent);}
        \end{center}
        \item \texttt{EmergencyMediator} đồng bộ thực thi 3 hành động an toàn khẩn cấp:
        \begin{itemize}
            \item Kích hoạt còi hú to cảnh báo nguy hiểm: \texttt{alarm.SoundAlarm(\dots)}.
            \item Tự động mở khóa toàn bộ cửa thoát hiểm: \texttt{doorController.UnlockAllDoors()}.
            \item Cách ly luồng gió, bật quạt hút khói sự cố: \texttt{ventilationController.IsolateVentilation()}.
        \end{itemize}
    \end{enumerate}
    \item \textbf{Ngược lại, khi khói chỉ ở mức \texttt{WARNING} ($50\% \le \text{khói} < 70\%$):} \texttt{EmergencyBridgeObserver} chỉ ghi nhận thông tin mà không kích hoạt Mediator, tránh báo động giả làm gián đoạn vận hành phòng máy.
\end{itemize}

\newpage

\section{Kết quả chạy và Kiểm thử hệ thống}

\subsection{Bảng tổng hợp kết quả kiểm thử tự động (T1 -- T8)}
Chương trình cài đặt bộ kiểm thử tự động toàn diện 8 ca kiểm thử (\textbf{T1 -- T8}) bao phủ 100\% các yêu cầu kỹ thuật và tiêu chí chấm điểm:

\begin{table}[H]
\centering
\footnotesize
\renewcommand{\arraystretch}{1.20}
\begin{tabular}{cp{3.4cm}p{4.6cm}p{4.6cm}c}
\toprule
\textbf{Mã} & \textbf{Tình huống kiểm thử} & \textbf{Kết quả mong đợi} & \textbf{Kết quả thực tế ghi nhận} & \textbf{Trạng thái} \\
\midrule
\textbf{T1} & Adapter: Quy đổi giá trị 0..1000 sang 0..100\% và đơn vị chuẩn & Đọc $735 \rightarrow 73.5$, đơn vị \texttt{"\% smoke"}, cập nhật động khi lớp cũ thay đổi & Giá trị đạt $73.5$, đơn vị chuẩn xác, đọc giá trị mới $92.0$ khi sửa lớp cũ & \textbf{PASS} \\
\midrule
\textbf{T2} & Ngưỡng an toàn: Mức khói $< 50\%$ là bình thường & Không phát tán \texttt{AlertEvent}, Dashboard không tăng số lần hiển thị & \texttt{DisplayCount = 0}, \texttt{LastAlert = null}, an toàn tuyệt đối & \textbf{PASS} \\
\midrule
\textbf{T3} & Ngưỡng cảnh báo: Khói từ $50\%$ đến dưới $70\%$ & Tạo sự kiện mức \texttt{"WARNING"}, gửi tới Dashboard & \texttt{LastAlert.Severity = "WARNING"}, giá trị $55.0\%$ smoke & \textbf{PASS} \\
\midrule
\textbf{T4} & Ngưỡng nguy cấp: Khói $\ge 70\%$ phát mức CRITICAL & Tạo sự kiện mức \texttt{"CRITICAL"}, gửi tới Dashboard & \texttt{LastAlert.Severity = "CRITICAL"}, giá trị $70.0\%$ smoke & \textbf{PASS} \\
\midrule
\textbf{T5} & Observer: Ngăn đăng ký trùng lặp và hỗ trợ Detach & Đăng ký cùng một observer 2 lần chỉ tính 1; Detach gỡ bỏ thành công & \texttt{ObserverCount = 2} khi đăng ký lặp; giảm về $1$ sau khi Detach & \textbf{PASS} \\
\midrule
\textbf{T6} & Observer Fault Isolation: Cách ly lỗi của một Observer & Observer lỗi ném ngoại lệ không làm sập vòng lặp; Observer khác vẫn nhận tin & Bắt ngoại lệ tại \texttt{FaultyObserver}; \texttt{goodObserver} vẫn nhận đủ sự kiện & \textbf{PASS} \\
\midrule
\textbf{T7} & Mediator: Điều phối đồng bộ 3 thiết bị khi có sự cố & Còi báo động hú, cửa mở khóa thoát nạn, quạt thông gió cách ly & Cả 3 cờ trạng thái \texttt{IsSounding}, \texttt{AreDoorsUnlocked}, \texttt{IsIsolated} đều \texttt{True} & \textbf{PASS} \\
\midrule
\textbf{T8} & Tích hợp E2E: Khói CRITICAL tự động kích hoạt Mediator & Khói WARNING không kích hoạt; khi tăng lên CRITICAL thì Bridge gọi Mediator & WARNING không gọi Mediator; khói $76\%$ kích hoạt đầy đủ 3 thiết bị & \textbf{PASS} \\
\bottomrule
\end{tabular}
\end{table}

\newpage

\subsection{Hình ảnh minh chứng thực nghiệm và Kiểm thử}

\subsubsection{1. Kịch bản mô phỏng vận hành thực tế (Phần Demo)}
Kịch bản mô phỏng trung tâm giám sát an toàn Campus Safety Hub: Tích hợp phần cứng cảm biến khói cũ qua Adapter, đăng ký các Observer, kiểm tra định kỳ mức khói an toàn, cảnh báo WARNING và tự động kích hoạt phản ứng khẩn cấp toàn diện khi đạt ngưỡng CRITICAL:

\begin{figure}[H]
\centering
\IfFileExists{demo.png}{%
    \includegraphics[width=0.92\textwidth]{demo.png}%
}{%
    \fbox{%
        \begin{minipage}[c][4.6cm][c]{0.88\textwidth}
            \centering
            \textbf{[VỊ TRÍ CHÈN ẢNH CHỤP MÀN HÌNH - PHẦN 1: KỊCH BẢN MÔ PHỎNG THỰC TẾ]}\\[0.2cm]
            \textit{(Tự động hiển thị hình ảnh khi lưu file \texttt{demo.png} cùng thư mục với file \texttt{BaoCao.tex})}\\[0.25cm]
            \small Thể hiện: Chuẩn hóa cảm biến qua Adapter, các tình huống khói an toàn, WARNING và kích hoạt hỏa hoạn CRITICAL điều phối qua Mediator.
        \end{minipage}%
    }%
}
\vspace{-0.2cm}
\caption{Kết quả thực thi kịch bản mô phỏng vận hành hệ thống Campus Safety Hub}
\label{fig:demo}
\end{figure}

\subsubsection{2. Kết quả thực thi bộ kiểm thử tự động (T1 -- T8)}
Chương trình thực thi bộ kiểm thử tự động đối soát tất cả các yêu cầu kỹ thuật và tiêu chí chấm điểm, đạt kết quả \textbf{PASS (8/8 ca kiểm thử)}:

\begin{figure}[H]
\centering
\IfFileExists{tests.png}{%
    \includegraphics[width=0.92\textwidth]{tests.png}%
}{%
    \fbox{%
        \begin{minipage}[c][4.8cm][c]{0.88\textwidth}
            \centering
            \textbf{[VỊ TRÍ CHÈN ẢNH CHỤP MÀN HÌNH - PHẦN 2: BỘ KIỂM THỬ TỰ ĐỘNG T1 - T8]}\\[0.2cm]
            \textit{(Tự động hiển thị hình ảnh khi lưu file \texttt{tests.png} cùng thư mục với file \texttt{BaoCao.tex})}\\[0.25cm]
            \small Thể hiện: Toàn bộ 8 ca kiểm thử T1 -- T8 vượt qua thành công với kết quả PASS màu xanh lá cây.
        \end{minipage}%
    }%
}
\vspace{-0.2cm}
\caption{Kết quả chạy bộ kiểm thử tự động đạt 8/8 ca kiểm thử chuẩn}
\label{fig:tests}
\end{figure}

\vspace{-0.1cm}

\section{Kết luận}
Bài thực hành Tuần 11 với hệ thống \textbf{Campus Safety Hub} đã ứng dụng xuất sắc 3 mẫu thiết kế phần mềm kinh điển nhằm giải quyết trọn vẹn bài toán tích hợp, giám sát và điều phối khẩn cấp:
\begin{itemize}[itemsep=1.0pt]
    \item \textbf{Adapter Pattern:} Cho phép tích hợp mượt mà cảm biến khói thế hệ cũ vào kiến trúc mới mà không làm sửa đổi hay lộ trạng thái nội bộ của phần cứng cũ.
    \item \textbf{Observer Pattern:} Phân tách triệt để bộ giám sát \texttt{SensorMonitor} khỏi các bên tiếp nhận thông báo, hỗ trợ đăng ký/hủy linh hoạt, cách ly lỗi hoàn hảo và bảo vệ dữ liệu với \texttt{AlertEvent} bất biến.
    \item \textbf{Mediator Pattern:} Loại bỏ cấu trúc phụ thuộc chéo "mạng nhện", tập trung hóa toàn bộ logic ứng phó khẩn cấp của phòng máy tại \texttt{EmergencyMediator}, bảo đảm còi báo động, cửa an toàn và hệ thống thông gió phối hợp nhịp nhàng, tin cậy.
\end{itemize}

\end{document}
